% Source: https://artofproblemsolving.com/community/c3h1485630_trapezium
% Stable identity: aops-topic-1485630-q01
A trapezoid inscribed in a circle is isosceles: the parallel chords $AB$ and $CD$ cut off equal arcs $AD$ and $BC$, so $AD=BC$.

\begin{asy}
usepackage("amsmath");
size(6.5cm);
defaultpen(fontsize(10pt));
path ra(pair P,pair u,pair v,real s=0.2){return P+s*unit(u)--P+s*unit(u)+s*unit(v)--P+s*unit(v);}
real R=8.2; real d1=2*aSin(8/(2*R)), d2=2*aSin(12/(2*R)), dpb=2*aSin(16/(2*R));
real t=270-(d2-d1)/2;
pair P=R*dir(t), D=R*dir(t-d1), C=R*dir(t+d2), B=R*dir(t+dpb), A=R*dir(540-(t+dpb));
pair M=extension(A,P,C,D);
draw(circle((0,0),R),gray);
draw(A--B--C--D--cycle,linewidth(0.9));
draw(A--P,blue); draw(P--B,red); draw(P--C,red); draw(P--D,red);
dot(A);dot(B);dot(C);dot(D);dot(P);dot(M,blue);
label("$A$",A,W); label("$B$",B,E); label("$C$",C,E); label("$D$",D,W); label("$P$",P,S); label("$M$",M,NE,blue);
label("$8$",D+0.3*(P-D),NE,red); label("$12$",(P+C)/2,SE,red); label("$16$",P+0.62*(B-P),SE,red);
label("\mbox{}",(min(currentpicture,user=true)+max(currentpicture,user=true))/2,black); // keeps the text after the figure black when LaTeX typesets the labels inline
\end{asy}

Compare triangles $PDM$ and $PBC$.
\begin{itemize}
\item $\angle DPM=\angle DPA$ and $\angle BPC$ are inscribed angles on the equal arcs $DA$ and $BC$ (the arcs not containing $P$), so $\angle DPM=\angle BPC$.
\item $\angle PDM=\angle PDC$ and $\angle PBC$ are inscribed angles on the same arc $PC$ (both $D$ and $B$ lie on the other side of chord $PC$), so $\angle PDM=\angle PBC$.
\end{itemize}
Hence triangles $PDM$ and $PBC$ are similar, and
\[
\dfrac{PM}{PC}=\dfrac{PD}{PB}\quad\Longrightarrow\quad PM=\dfrac{PC\cdot PD}{PB}=\dfrac{12\cdot8}{16}=\boxed{6}.
\]
